\section{Control Unit}

\vfill
\begin{figure}[H]
  \centering
  \begin{adjustbox}{max width=\textwidth}
  \begin{tikzpicture}[processor arch]
    \node[register, anchor=north east, minimum width=15cm] (datapath) {DataPath};

    \node[register, anchor=south west, yshift=4cm, minimum width=4.5cm, minimum height=3cm] at (datapath.west |- datapath.north) (i_decoder) {Instruction\\ Decoder};
    \node[register, anchor=south west, yshift=1.5cm] at (i_decoder.north west) (step_counter) {Step\\ Couter};

    \node[register, anchor=south, yshift=4cm] at (datapath.north) (mux_pc) {MUX\\ PC};
    \node[register] at (mux_pc.south |- step_counter.east) (pc) {PC};

    \begin{registergroup}{ir}{IR (64)}
      \node[register, anchor=west, yshift=7cm] at (step_counter.west) (ir_1) {IR\\ {[Op]}};
      \node[register, right=5pt of ir_1] (ir_2) {IR\\ {[Arg1]}};
      \node[register, right=5pt of ir_2] (ir_3) {IR\\ {[Arg1/2]}};
      \node[register, right=5pt of ir_3] (ir_4) {IR\\ {[-/Arg2]}};
    \end{registergroup}

    \node[register, anchor=south, yshift=1.5cm] at (ir.north) (mux_ir) {MUX\\ IR};

    \node[memory, minimum height=6cm, anchor=south east, yshift=12cm] at (datapath.north east) (mem) {Memory:\\ \textit{Instructions,}\\ \textit{Data}};
  \end{tikzpicture}
  \end{adjustbox}
\end{figure}
\vfill

\newpage

* счётчик тактов (для инструкций, которые выполняются за несколько тактов)
* это у меня с состоянием или без него? (в данном процессоре)

Комбинационная схема для CISC (т.к. Пенской считает, что микрокодом будет долго)

Непрямое отображение кода операции на адрес микроинструкции (см. 8:12-5)